\documentclass{article}
\usepackage{../mathnotezh}

\title{空间直线夹角的转化}
\author{tbj}
\date{}

\begin{document}

\maketitle

对于空间中四点$A, B, C, D$,要想求直线$AB, CD$的夹角,只要知道四个点两两之间的距离即可.
具体地说,我们有如下定理

\begin{thm}
    对于空间中四点$A, B, C, D$,设直线$AB, CD$之间的夹角为$\theta$,则
    \begin{equation}
        \cos \theta = \dfrac{\lrvert{\lrvert{AD}^2 + \lrvert{BC}^2 -
        \lrvert{AC}^2 - \lrvert{BD}^2}}{2 \lrvert{AB} \lrvert{CD}} \label{1}
    \end{equation}
\end{thm}

我们用三种方法证明该定理.

\paragraph{几何法}

将$CD$平移至$EF$,使得$A, B, E, F$四点共面,且$CE, DF$垂直于该平面.
设直线$AB, EF$交于点$P$,则$\cos \theta = \lrvert{\cos P}$,
且$\lrvert{CD} = \lrvert{EF}, \lrvert{CE} = \lrvert{DF} = d$.

\begin{figure}[h]
    \centering
    \includegraphics[width=0.7\textwidth]{fig1.pdf}
\end{figure}

因此我们有
\begin{align*}
    \lrvert{AC}^2 & = \lrvert{AE}^2 + d^2 \\
    \lrvert{AD}^2 & = \lrvert{AF}^2 + d^2 \\
    \lrvert{BC}^2 & = \lrvert{BE}^2 + d^2 \\
    \lrvert{BD}^2 & = \lrvert{BF}^2 + d^2 \\
    \lrvert{AE}^2 & = \lrvert{PA}^2 + \lrvert{PE}^2 - 2 \lrvert{PA} \lrvert{PE} \cos P \\
    \lrvert{AF}^2 & = \lrvert{PA}^2 + \lrvert{PF}^2 - 2 \lrvert{PA} \lrvert{PF} \cos P \\
    \lrvert{BE}^2 & = \lrvert{PB}^2 + \lrvert{PE}^2 - 2 \lrvert{PB} \lrvert{PE} \cos P \\
    \lrvert{BF}^2 & = \lrvert{PB}^2 + \lrvert{PF}^2 - 2 \lrvert{PB} \lrvert{PF} \cos P
\end{align*}

进而
\begin{align*}
    & \dfrac{\lrvert{\lrvert{AD}^2 + \lrvert{BC}^2 - \lrvert{AC}^2 - \lrvert{BD}^2}}{2 \lrvert{AB} \lrvert{CD}} \\
    & = \dfrac{\lrvert{\lrvert{AF}^2 + \lrvert{BE}^2 - \lrvert{AE}^2 - \lrvert{BF}^2}}{2 \lrvert{AB} \lrvert{CD}} \\
    & = \lrvert{\cos P} \dfrac{\lrvert{\lrvert{PA} \lrvert{PE} + \lrvert{PB} \lrvert{PF} -
        \lrvert{PA} \lrvert{PF} - \lrvert{PB} \lrvert{PE}}}{\lrvert{AB} \lrvert{EF}} \\
    & = \lrvert{\cos P} \dfrac{\lrvert{\lrvert{PA} - \lrvert{PB}} \lrvert{\lrvert{PE} -
        \lrvert{PF}}}{\lrvert{AB} \lrvert{EF}} \\
    & = \lrvert{\cos P} \\
    & = \cos \theta \tag*{\qed}
\end{align*}

\begin{rqq}
    上述方法运用了勾股定理,余弦定理,计算繁复,接下来给出简单一点的方法.
\end{rqq}

\paragraph{向量法1}

事实上,可以证明
\begin{equation}
    2 \vec{AB} \cdot \vec{CD} = \vec{AD}^2 + \vec{BC}^2 -
    \vec{AC}^2 - \vec{BD}^2 \label{2}
\end{equation}

我们有
\begin{align*}
    & \vec{AD}^2 + \vec{BC}^2 - \vec{AC}^2 - \vec{BD}^2 \\
    & = \left( \vec{AD}^2 - \vec{AC}^2 \right) - \left( \vec{BD}^2 - \vec{BC}^2 \right) \\
    & = \left( \vec{AD} + \vec{AC} \right) \left( \vec{AD} - \vec{AC} \right) -
        \left( \vec{BD} + \vec{BC} \right) \left( \vec{BD} - \vec{BC} \right) \\
    & = \left( (\vec{AC} - \vec{BC}) + (\vec{AD} - \vec{BD}) \right)
        \cdot \vec{CD} \\
    & = 2 \vec{AB} \cdot \vec{CD}
\end{align*}

\eqref{2} 得证.两边同除$2 \lrvert{\vec{AB}} \lrvert{\vec{CD}}$并注意到
$\cos \theta = \lrvert{\cos \langle \vec{AB}, \vec{CD} \rangle}$即证 \eqref{1}. \hfill\qed

\paragraph{向量法2}

设$O$为原点(事实上可以为空间中任意一点),要证 \eqref{2},只要证
\[ 2(\vec{OB} - \vec{OA})(\vec{OD} - \vec{OC}) = (\vec{OD} - \vec{OA})^2 + (\vec{OC} - \vec{OB})^2 -
(\vec{OC} - \vec{OA})^2 - (\vec{OD} - \vec{OB})^2 \]

展开即得. \hfill\qed

\begin{rqq}
    上述结论的一种特殊情况是$AB \perp CD \iff \lrvert{AC}^2 - \lrvert{AD}^2 = \lrvert{BC}^2 - \lrvert{BD}^2$,
    这在平面几何中称为\textbf{定差幂线引理}.
\end{rqq}

\begin{rqq}
    上述结论的另一种特殊情况是$AB \parallel CD \iff \lrvert{\lrvert{AD}^2 + \lrvert{BC}^2 - \lrvert{AC}^2 - \lrvert{BD}^2} =
    2\lrvert{AB} \lrvert{CD}$.
\end{rqq}

\begin{rqq}
    点$A, B, C, D$可以在直线上任意选取(不能重合),不改变两直线的夹角.这意味着,
    $\dfrac{\lrvert{\lrvert{AD}^2 + \lrvert{BC}^2 - \lrvert{AC}^2 - \lrvert{BD}^2}}{2 \lrvert{AB} \lrvert{CD}}$
    为一种不变量.
\end{rqq}

\begin{rqq}
    向量法2对 \eqref{2} 的证明没有涉及三维空间的性质,也没有涉及空间中的作图,这表明我们的结论在高维空间中仍然成立.
\end{rqq}

\end{document}